%%%%%%%%%%%%%%%%%%%%%%%%%%%%%%%%%%%%%%%%%%%%%%%%%%%%%%%%%%%%%%%%%%%%%%
% Lớp 8 | Chương I. Đa thức | Phép cộng và phép trừ đa thức
% Dạng 8D1?3-3: Bài toán thực tế dùng cộng, trừ đa thức
%%%%%%%%%%%%%%%%%%%%%%%%%%%%%%%%%%%%%%%%%%%%%%%%%%%%%%%%%%%%%%%%%%%%%%


% nguồn: De_AI_Toan8_cuoi_hoc_ki_i_1 (nhập 21/11/2026)
\begin{bt}%[8D1H3-3]%[ma:73c6dd22]
	Một khu vườn hình tam giác có độ dài ba cạnh lần lượt là $2x + y$, $x + 3y$ và $3x - y$ (mét).
a) Viết đa thức thu gọn biểu thị chu vi của khu vườn.
b) Tính chu vi khu vườn khi $x = 4$, $y = 2$.
	\loigiai{
		a) Chu vi khu vườn là $C = (2x + y) + (x + 3y) + (3x - y) = 6x + 3y$ (m). \\
b) Thay $x = 4$, $y = 2$ vào $C$, ta được $C = 6 \cdot 4 + 3 \cdot 2 = 24 + 6 = 30\text{ m}$.
	}
\end{bt}

% nguồn: De_AI_Toan8_cuoi_hoc_ki_i_1 (nhập 21/11/2026)
% nguồn: De_AI_Toan8_cuoi_hoc_ki_i_1 (nhập 21/11/2026)
\begin{bt}%[8D1H3-3]%[ma:b25ee303]
	Một hình hộp chữ nhật có thể tích là $6x^3y + 9x^2y^2$ ($\text{cm}^3$) và diện tích mặt đáy là $3x^2y$ ($\text{cm}^2$) với $x > 0$, $y > 0$.
a) Viết đa thức biểu thị chiều cao của hình hộp chữ nhật.
b) Tính chiều cao của hình hộp khi $x = 2$, $y = 1$.
	\loigiai{
		a) Chiều cao của hình hộp chữ nhật là $h = (6x^3y + 9x^2y^2) : (3x^2y) = 2x + 3y$ (cm). \\
b) Thay $x = 2$, $y = 1$ vào $h$, ta được $h = 2 \cdot 2 + 3 \cdot 1 = 7\text{ cm}$.
	}
\end{bt}

% nguồn: De_AI_Toan8_cuoi_hoc_ki_i_1 (nhập 21/11/2026)
\begin{bt}%[8D1H3-3]%[ma:ba33f5cf]
	Một chiếc bình dạng hình lập phương có cạnh $2x$ (cm). Mực nước trong bình đang cao $y$ (cm) với $0 < y < 2x$.
a) Viết đa thức biểu thị thể tích phần bình chưa có nước.
b) Tính thể tích phần bình chưa có nước khi $x = 5$, $y = 4$.
	\loigiai{
		a) Thể tích cả bình là $(2x)^3 = 8x^3$ ($\text{cm}^3$). Thể tích nước trong bình là $(2x)^2 \cdot y = 4x^2y$ ($\text{cm}^3$). \\
Thể tích phần bình chưa có nước là $V = 8x^3 - 4x^2y$ ($\text{cm}^3$). \\
b) Thay $x = 5$, $y = 4$ vào $V$, ta được $V = 8 \cdot 5^3 - 4 \cdot 5^2 \cdot 4 = 1000 - 400 = 600\text{ cm}^3$.
	}
\end{bt}

% nguồn: De_AI_Toan8_cuoi_hoc_ki_i_1 (nhập 21/11/2026)
\begin{bt}%[8D1H3-3]%[ma:82862cb0]
	Một cửa hàng bán mỗi hộp bánh giá $50$ nghìn đồng và mỗi hộp sữa giá $30$ nghìn đồng. Ngày thứ nhất cửa hàng bán được $x$ hộp bánh và $y$ hộp sữa; ngày thứ hai bán được $x + 5$ hộp bánh và $2y - 3$ hộp sữa ($y \ge 2$).
a) Viết đa thức thu gọn biểu thị tổng doanh thu (nghìn đồng) của cả hai ngày.
b) Tính tổng doanh thu của cả hai ngày khi $x = 10$, $y = 5$.
	\loigiai{
		a) Doanh thu ngày thứ nhất là $50x + 30y$ (nghìn đồng). \\
Doanh thu ngày thứ hai là $50(x + 5) + 30(2y - 3) = 50x + 250 + 60y - 90 = 50x + 60y + 160$ (nghìn đồng). \\
Tổng doanh thu hai ngày là $(50x + 30y) + (50x + 60y + 160) = 100x + 90y + 160$ (nghìn đồng). \\
b) Thay $x = 10$, $y = 5$ vào biểu thức trên, ta được $100 \cdot 10 + 90 \cdot 5 + 160 = 1610$ (nghìn đồng).
	}
\end{bt}
